\scanpage{106}
% archive foliation: 105
Donc \deleted{\unsure{négligeant}} pour $p+q=i-1$, on trouve un poids $\le2(n-1)-i+1=2n-i-1$.

\uline{Remarque} La filtration de $H^{\unsure{*}}(Y)$ déduite de la suite spectrale \unsure{est} n'est autre que la filtration par les poids. La filtration image de $H^i_!$ est donc aussi la filtration par les poids croissants (de $0$ à $i$, le dernier avec comme \unsure{noyau}) donc est également indépendante de la \unsure{comp.} \unsure{projective} $X'$ choisie. Écrivons les termes de la filtration, en utilisant le fait que $H^{i-1}(X')$ est purement de poids $i-1$, et $\deleted{\unsure{E^{i-1,0}_2}}$ que la suite spectrale est dégénérée par des raisons de poids :

\diagram{\(\begin{array}{c|c|c|c|c|c|c}
\hline
& E^{i-1,0}_2, & E^{i-2,1}_2, & \ldots & E^{1,i-2}_2, & E^{0,i-1}_2/\operatorname{Im}H^{i-1}(X'), & \operatorname{Ker}(H^i(X')\to E^{0,i}_2)\\
\hline
\text{poids} & 0 & 1 & & i-2 & i-1 & i\\
\hline\hline
\deleted{\unsure{\text{niveau}}} & & & & & &
\end{array}\)}

\deleted{\unsure{Comme $i\ge n$}} Si $i\le n$, les niveaux sont \inserted{seulement} $\le$ poids, il n'y a pas lieu à priori d'attendre plus. L'on dresse d'autre part, pour $i\ge n$, un tableau

\diagram{\(\begin{array}{c|c|c|c|c|c|c}
& \text{degré }\unsure{\text{filtrant}}\ge n & & & & &\\
& 0\ \ldots\ 0 &
\substack{\deleted{\unsure{E^{i-1,0}_2}}\\E^{2n-i-1,2(i-n)}_2} &
E^{2n-i-2,2(i-n)+1}_2 &
\ldots &
E^{0,i-1}_2/\operatorname{Im}H^{i-1}(X') &
\operatorname{Ker}(H^i(X')\to E^{0,i}_2)\\
\hline
\text{poids} & \raisebox{-1.8ex}{\Large$\times$} & 2(i-n) & 2(i-n)+1 & \ldots & i-1 & i\\
\cline{3-7}
\text{niveau} & & 0 & 1 & \ldots & 2n-i-1 & 2n-i
\end{array}\)}

\uline{Corollaire 1} (Par dualité) $H^i(X)$ est de poids compris entre $i$ et $2n$. Le morceau \deleted{\unsure{pur}} de poids $i$ est exactement l'image de $H^i(X')$. Si $i\unsure{\ge|>}n$, les niveaux des morceaux de poids $i,i+1,\ldots,2n$ sont resp. $\le2n-i,2n-i-1,\ldots,0$. Si $i\le n$ il n'y a de morceaux \deleted{\unsure{que}} que des poids
$\deleted{i,i+1,\ldots,2n\ \unsure{2(i-n)=2n-2n}}$ \deleted{\unsure{dont}}
\inserted{$i,i+1,\ldots,2i$}
qui sont de niveaux resp. $\le i,i-1,\ldots,0$.

\margin{\fbox{\parbox{.24\linewidth}{tous les morceaux de la filtration \unsure{st} ``effectifs''}}}
\margin{\fbox{\parbox{.24\linewidth}{En tous cas niveau $\le\operatorname{Inf}(i,n)$}}}
\margin{\uline{NB} les niveaux ici \uline{décroissent}}
